\section{Intersections of spherical caps}
\label{sec:caps}

We need two properties of the cutting caps: separated caps have small
overlap, whereas a collection of nearby caps has a substantial common
intersection. For a target $y\in G_P$ and a relevant center $p$, write
$x=y-p$. The set
\[
 C(x)=\{u\in S^{n-1}:|u\cdot x|>1+3\eta\}
\]
consists of the normals whose slabs exclude $y$ from $p+K^+$.
Figure~\ref{fig:residual-cap} illustrates this correspondence. The
estimates below allow the slightly more general thresholds needed to
control these caps uniformly.
\begin{figure}[htbp]
\centering
\begin{tikzpicture}[x=1cm,y=1cm,>=Stealth,
 every node/.style={font=\small},
 axis/.style={draw=black!45,line width=.5pt},
 vector/.style={->,draw=blue!55!black,line width=.9pt}]
 \coordinate (p) at (0,0);
 \coordinate (y) at (2.9,.65);
 \fill (p) circle (1.5pt) node[below left] {$p$};
 \fill (y) circle (1.5pt) node[above right] {$y$};
 \draw[vector] (p) -- (y) node[midway,above,sloped] {$x=y-p$};
 \node[align=center] at (1.45,-1.05) {center and target\\in physical space};

 \begin{scope}[shift={(7.9,0)}]
  \fill[red!18] (35:1.4) arc (35:-35:1.4) -- cycle;
  \fill[red!18] (145:1.4) arc (145:215:1.4) -- cycle;
  \draw[axis] (0,0) circle (1.4);
  \draw[axis,->] (-1.8,0) -- (1.95,0) node[right] {$x/|x|$};
  \draw[dashed,black!55] (35:1.4) -- (-35:1.4);
  \draw[dashed,black!55] (145:1.4) -- (215:1.4);
  \draw[vector] (0,0) -- (20:1.4) node[above right] {$u$};
  \node[align=center] at (0,-2.05)
    {cutting normals $C(x)$\\$|u\cdot x|>1+3\eta$};
  \node at (0,1.8) {normal sphere};
 \end{scope}
\end{tikzpicture}
\caption{Plane-section schematic of a residual and its cutting cap.
The shaded antipodal caps lie beyond the two threshold hyperplanes
$u\cdot x=\pm(1+3\eta)$. A sampled normal in $C(x)$ excludes the target
$y$ from $p+K^+$. The drawing does not depict high-dimensional cap
measure. For a block of residuals, one normal in their common subset
cuts them all; assigning distinct normals to all blocks is a sufficient
certificate of a hole.}
\label{fig:residual-cap}
\end{figure}


The estimates adapt Li--Liu's cap and common-lens
arguments \cite[Lemma~6.1 and Lemmas~9.2--9.3]{LiLiu2026Lattice}.
In our thin radial shell, a bisector argument gives a short proof of the
overlap estimate with a constant prefactor. Keeping that prefactor
independent of dimension will be important when we sum overlaps.

\begin{lemma}[Cap estimates]\label{lem:caps}
Fix \(1<a<2\) and \(A,Q,B_0>0\). There exist constants
\(C_a,c_a'>0\), \(C_{a,A}>0\), and \(L=L(a)>0\) such that the
following holds for all sufficiently large \(n\), depending on
\(a,A,Q,B_0\). Let \(\sigma\) be normalized surface measure on
\(S^{n-1}\), and set
\[
 M=\sigma\{u:|u\cdot e_1|>1/a\}^{-1},\qquad \mu=M\sigma.
\]
Let \(x_1,\ldots,x_m\in\mathbb{R}^n\) and
\(\epsilon_1,\ldots,\epsilon_m\in\mathbb{R}\), where
\(1\le m\le n^2\), satisfy
\[
 a\left(1-\frac{A\log n}{n}\right)-Qn^{-4}
 \le |x_i|\le a(1+1/n)+Qn^{-4},\qquad
 |\epsilon_i|\le Qn^{-4}.
\]
Define
\[
 C_i=\{u\in S^{n-1}:|u\cdot x_i|>1+\epsilon_i\},\quad
 c_i=\mu(C_i),\quad \bar c_i=\min(c_i,1/2),\quad
 s_i=-\log\bar c_i,
\]
and, using projective angles, define
\[
 d(i,j)=\sqrt n\,\angle(x_i,x_j)+\sqrt{|s_i-s_j|},\qquad
 r(S)=\min_{k\in S}\max_{i\in S}d(k,i)^2
\]
for each nonempty set \(S\) of indices. Then
\begin{align}
 n^{-C_{a,A}}\le c_i&\le C_a,
 &s_i&\le C_a\bigl(1+n[1-|x_i|^2/a^2]_+\bigr),
 \label{eq:caps-cost}\\
 \frac{\mu(C_i\cap C_j)}{\sqrt{\bar c_i\bar c_j}}
 &\le C_a\exp\{-c_a'd(i,j)^2\}.
 \label{eq:caps-affinity}
\end{align}
The function \(d\) is a pseudometric on the indexed list. If
\(r(S)\le B_0\log n\) and \(i_0\in S\) minimizes the raw intensity
\(c_i\), then
\begin{equation}
 \mu\left(\bigcap_{i\in S}C_i\right)
 \ge e^{-L(1+r(S)+\log|S|)}\bar c_{i_0}.
 \label{eq:caps-lens}
\end{equation}
In particular, \(L\) is independent of \(A,Q,B_0\); these parameters
may affect the starting dimension.
\end{lemma}

\begin{proof}
We first record one-dimensional estimates, then prove the overlap and
common-intersection bounds. Constants with a subscript \(a\) below may
change from line to line, but depend only on \(a\).

\paragraph{One-cap probabilities and radial costs.}
Write \(t_*=1/a\), \(e_i=x_i/|x_i|\), and
\(\tau_i=(1+\epsilon_i)/|x_i|\). For a unit vector \(e\), let
\[
 F_n(t)=\sigma\{u:u\cdot e>t\}
 =k_n\int_t^1(1-v^2)^{(n-3)/2}\,dv,\qquad 0<t<1,
\]
where \(k_n\) is the normalization constant. Substitution of
\(z=1-v^2\), using \(1\le1/v\le1/t\), gives
\begin{equation}
 \frac{k_n}{n-1}(1-t^2)^{(n-1)/2}
 \le F_n(t)\le
 \frac{k_n}{(n-1)t}(1-t^2)^{(n-1)/2}.
 \label{eq:caps-mills}
\end{equation}
Consequently,
\begin{equation}
 \frac{(n-1)t}{1-t^2}
 \le-\frac{d}{dt}\log F_n(t)
 \le\frac{n-1}{1-t^2}.
 \label{eq:caps-hazard}
\end{equation}
These exact inequalities avoid dimension-dependent prefactors in
ratios of cap probabilities.

All the thresholds satisfy
\begin{equation}
 \tau_i=t_*+O_{a,A,Q}(\log n/n),\qquad
 \tau_i\ge t_*-C_a/n,
 \label{eq:caps-thresholds}
\end{equation}
for sufficiently large \(n\), and hence belong to the fixed interval
\(J=[t_*/2,(1+t_*)/2]\subset(0,1)\).
Since \(M=(2F_n(t_*))^{-1}\), we have
\(c_i=F_n(\tau_i)/F_n(t_*)\).
The logarithmic derivative in \eqref{eq:caps-hazard} is bounded by
\(C_an\) on \(J\). This proves \(c_i\le C_a\), while the lower radial
bound proves \(c_i\ge n^{-C_{a,A}}\).
More explicitly, putting \(\delta_i=[1-|x_i|^2/a^2]_+\), elementary
algebra gives
\[
\begin{gathered}
 (\tau_i-t_*)_+\le C_a(\delta_i+Qn^{-4}),\qquad
 (t_*-\tau_i)_+\le C_a(n^{-1}+Qn^{-4}),\\
 n\delta_i\le C_{a,A}\log n+O_{a,Q}(n^{-3}).
\end{gathered}
\]
Integrating \eqref{eq:caps-hazard} therefore gives
\(-\log c_i\le C_a(1+n\delta_i)\), after increasing the starting
dimension so that \(Qn^{-3}\le1\). As
\(s_i=\max(\log2,-\log c_i)\), this proves
\eqref{eq:caps-cost}. The additive constant is needed even when
\(c_i\ge1/2\).

\paragraph{The overlap bound.}
Consider two ordinary half caps with thresholds \(\tau,\upsilon\)
drawn from our list and with axis angle \(\theta\in[0,\pi]\).
Put \(t=(\tau+\upsilon)/2\). If \(\theta<\pi\), adding their two
defining inequalities shows that their intersection lies in the half
cap along the angle bisector with threshold
\[
 q=t/\cos(\theta/2).
\]
The intersection is empty if \(\theta=\pi\) or \(q\ge1\), so suppose
\(q<1\). Let \(f(z)=\log(1-z^2)\). Its second derivative is bounded
in absolute value on \(J\), whence
\[
 f(t)-\tfrac12 f(\tau)-\tfrac12 f(\upsilon)
 \le C_a(\tau-\upsilon)^2.
\]
Furthermore,
\[
 \frac{1-q^2}{1-t^2}
 =1-\frac{t^2}{1-t^2}\tan^2(\theta/2)
 \le e^{-c_a'\theta^2},
\]
using \(t\ge t_*/2\) and \(\tan(\theta/2)\ge\theta/2\).
Apply the upper bound in \eqref{eq:caps-mills} to \(F_n(q)\) and
the lower bound to both terms in the following denominator:
\begin{equation}
 \frac{F_n(q)}{\sqrt{F_n(\tau)F_n(\upsilon)}}
 \le q^{-1}\exp\{C_an(\tau-\upsilon)^2-c_a'n\theta^2\}
 \le C_a e^{-c_a'n\theta^2}.
 \label{eq:caps-bisector}
\end{equation}
Indeed, \(q^{-1}\le2/t_*\), and
\(n(\tau-\upsilon)^2=O_{a,A,Q}((\log n)^2/n)\to0\) by
\eqref{eq:caps-thresholds}. This last fact is the benefit of the
thin shell. The argument does not require \(q\in J\), so it also
applies when the bisector threshold approaches 1.

For projective angle \(\phi=\angle(x_i,x_j)\), the four pairs of
signed half caps have ordinary angles \(\phi\) or
\(\pi-\phi\ge\phi\). Summing \eqref{eq:caps-bisector} over these
four disjoint sign choices and using \(c_i=2MF_n(\tau_i)\) gives
\[
 \mu(C_i\cap C_j)\le C_a\sqrt{c_ic_j}\,e^{-c_a'n\phi^2}.
\]
By \eqref{eq:caps-cost}, \(c_i/\bar c_i\le\max(1,2C_a)\), so
clipping the denominator costs only a constant. Independently, the
trivial intersection bound gives
\[
 \frac{\mu(C_i\cap C_j)}{\sqrt{\bar c_i\bar c_j}}
 \le C_a e^{-|s_i-s_j|/2}.
\]
To see this, order the clipped intensities as
\(\bar c_i\le\bar c_j\). If \(\bar c_i<1/2\), use
\(\mu(C_i\cap C_j)\le c_i=\bar c_i\); otherwise both clipped
intensities equal \(1/2\), and use \(c_i\le C_a\).
The minimum of the two bounds is at most their geometric mean.
Together with
\[
 d(i,j)^2\le2\bigl(n\phi^2+|s_i-s_j|\bigr),
\]
this proves \eqref{eq:caps-affinity}.

Projective angle satisfies the triangle inequality, as the quotient
of spherical distance under antipodal identification. The inequality
\(\sqrt{x+y}\le\sqrt x+\sqrt y\) shows that
\(\sqrt{|s_i-s_j|}\) also satisfies it. Thus \(d\) is a
pseudometric; distinct labels, including repeated caps, may have
distance zero.

\paragraph{A common intersection.}
Set \(r=r(S)\) and \(m_S=|S|\), and choose an index attaining the
minimum defining \(r\). By the triangle inequality, each axis in
\(S\) has projective angle at most \(2\sqrt{r/n}\) from
\(e_0=e_{i_0}\). Since \(F_n\) is strictly decreasing, the
raw-minimum anchor has the largest threshold:
\(\tau_0=\tau_{i_0}\ge\tau_i\).
Orient each axis so that its ordinary angle \(\phi_i\) with
\(e_0\) equals its projective angle, and write
\[
 e_i=(\cos\phi_i)e_0+(\sin\phi_i)v_i,\qquad
 \Phi=\max_{i\in S}\phi_i\le2\sqrt{r/n},
\]
where \(v_i\) is a unit tangent vector. For \(\phi_i=0\), its
choice is immaterial.

We use the elementary bound
\begin{equation}
 \mathbb{P}(W\cdot v<-h)\le e^{-dh^2/2}
 \label{eq:caps-tangent}
\end{equation}
for a uniform \(W\in S^{d-1}\), \(d\ge2\), any unit \(v\), and
\(h\ge0\). Here is a proof. Rotational invariance reduces to
\(X=W_1\). Integrating the derivative of
\(x^{2j-1}(1-x^2)^{(d-1)/2}\) on \([-1,1]\) shows that
\[
 \mathbb{E} X^{2j}
 =\frac{2j-1}{d+2j-2}\mathbb{E} X^{2j-2}
 \le\frac{(2j)!}{2^j j!d^j}.
\]
Odd moments vanish. The exponential series therefore gives
\(\mathbb{E} e^{zX}\le e^{z^2/(2d)}\); applying Markov's inequality
with \(z=-dh\) proves \eqref{eq:caps-tangent}.

Conditionally on \(u\cdot e_0=t\), a uniform sphere point has the
form \(u=te_0+\sqrt{1-t^2}\,W\), with \(W\) uniform on the tangent
sphere \(S^{n-2}\). Put
\[
 h=\sqrt{\frac{2\log(2m_S)}{n-1}},\qquad
 \delta=2\Phi^2+4\Phi h.
\]
By \eqref{eq:caps-tangent} and a union bound, with conditional
probability at least \(1/2\) one has \(W\cdot v_i\ge-h\) for
every \(i\in S\). For large enough \(n\),
\(\cos\Phi\ge1/2\) and \(\tau_0+\delta\in J\), because
\(r\le B_0\log n\) and \(m_S\le n^2\). On that tangent event,
whenever \(t>\tau_0+\delta\),
\[
 u\cdot e_i
 \ge t\cos\phi_i-h\sin\phi_i
 >\tau_0-\Phi^2/2+\delta/2-h\Phi
 \ge\tau_0\ge\tau_i.
\]
Thus all these points lie in every double cap. Integration over the
axial coordinate gives
\[
 \mu\left(\bigcap_{i\in S}C_i\right)
 \ge\frac M2F_n(\tau_0+\delta)
 =\frac{c_{i_0}}4\frac{F_n(\tau_0+\delta)}{F_n(\tau_0)}
 \ge\frac{c_{i_0}}4 e^{-C_an\delta},
\]
where the last step follows from \eqref{eq:caps-hazard} on \(J\).
Finally,
\[
 n\delta\le C\bigl(r+\sqrt{r\log(2m_S)}\bigr)
 \le C\bigl(r+\log(2m_S)\bigr).
\]
Absorb the fixed factors into \(L=L(a)\) and use
\(c_{i_0}\ge\bar c_{i_0}\) to obtain \eqref{eq:caps-lens}.
If \(r=0\), then \(\Phi=\delta=0\), and the same proof applies.
\end{proof}

For the parameters of the construction, the overlap vanishes beyond
a fixed angle. We record this separately because zero overlap gives
more than the exponentially small bound in Lemma~\ref{lem:caps}.

\begin{corollary}\label{cor:caps-disjoint}
With \(\varepsilon=2^{-30}\) and \(a=1+\varepsilon\), caps satisfying
the hypotheses of Lemma~\ref{lem:caps} are disjoint whenever their
axes have projective angle greater than \(4\sqrt\varepsilon\), for
all sufficiently large \(n\).
\end{corollary}

\begin{proof}
Each half cap has angular radius \(\arccos\tau_i\). Since
\[
 \cos(2\sqrt\varepsilon)
 \le1-2\varepsilon+\tfrac23\varepsilon^2
 <1-\varepsilon\le\frac1{1+\varepsilon},
\]
the estimate \(\tau_i\ge1/a-C_a/n\) shows that its radius is
strictly less than \(2\sqrt\varepsilon\) for large \(n\).
Every pair of signed axes is at ordinary angle at least the
projective angle. The spherical triangle inequality therefore
excludes a point in any pair of half caps at projective separation
greater than \(4\sqrt\varepsilon\).
\end{proof}
